\documentclass[10pt,a4paper]{amsart}
\usepackage{amsmath,amssymb,amsthm,graphicx}
% For commutative diagrams we use Xy-pic
\usepackage[all]{xy}

% We use 2cell for 2-commutative diagrams.
\xyoption{2cell}
% Theorem environments.
%
\theoremstyle{plain}
\newtheorem{theorem}[subsection]{Theorem}
\newtheorem{proposition}[subsection]{Proposition}
\newtheorem{lemma}[subsection]{Lemma}

\theoremstyle{definition}
\newtheorem{definition}[subsection]{Definition}
\newtheorem{example}[subsection]{Example}
\newtheorem{exercise}[subsection]{Exercise}
\newtheorem{situation}[subsection]{Situation}

\theoremstyle{remark}
\newtheorem{remark}[subsection]{Remark}
\newtheorem{remarks}[subsection]{Remarks}

\numberwithin{equation}{subsection}

% Macros
%
\def\lim{\mathop{\mathrm{lim}}\nolimits}
\def\colim{\mathop{\mathrm{colim}}\nolimits}
\def\Spec{\mathop{\mathrm{Spec}}}
\def\Hom{\mathop{\mathrm{Hom}}\nolimits}
\def\Ext{\mathop{\mathrm{Ext}}\nolimits}
\def\SheafHom{\mathop{\mathcal{H}\!\mathit{om}}\nolimits}
\def\SheafExt{\mathop{\mathcal{E}\!\mathit{xt}}\nolimits}
\def\Sch{\mathit{Sch}}
\def\Mor{\mathop{\mathrm{Mor}}\nolimits}
\def\Ob{\mathop{\mathrm{Ob}}\nolimits}
\def\Sh{\mathop{\mathit{Sh}}\nolimits}
\def\NL{\mathop{N\!L}\nolimits}
\def\CH{\mathop{\mathrm{CH}}\nolimits}
\def\proetale{{pro\text{-}\acute{e}tale}}
\def\etale{{\acute{e}tale}}
\def\QCoh{\mathit{QCoh}}
\def\Ker{\mathop{\mathrm{Ker}}}
\def\Im{\mathop{\mathrm{Im}}}
\def\Coker{\mathop{\mathrm{Coker}}}
\def\Coim{\mathop{\mathrm{Coim}}}

% Boxtimes
%
\DeclareMathSymbol{\boxtimes}{\mathbin}{AMSa}{"02}

%
% Macros for moduli stacks/spaces
%
\def\QCohstack{\mathcal{QC}\!\mathit{oh}}
\def\Cohstack{\mathcal{C}\!\mathit{oh}}
\def\Spacesstack{\mathcal{S}\!\mathit{paces}}
\def\Quotfunctor{\mathrm{Quot}}
\def\Hilbfunctor{\mathrm{Hilb}}
\def\Curvesstack{\mathcal{C}\!\mathit{urves}}
\def\Polarizedstack{\mathcal{P}\!\mathit{olarized}}
\def\Complexesstack{\mathcal{C}\!\mathit{omplexes}}
% \Pic is the operator that assigns to X its picard group, usage \Pic(X)
% \Picardstack_{X/B} denotes the Picard stack of X over B
% \Picardfunctor_{X/B} denotes the Picard functor of X over B
\def\Pic{\mathop{\mathrm{Pic}}\nolimits}
\def\Picardstack{\mathcal{P}\!\mathit{ic}}
\def\Picardfunctor{\mathrm{Pic}}
\def\Deformationcategory{\mathcal{D}\!\mathit{ef}}
\newenvironment{reference}{}{}
\newenvironment{slogan}{}{}
\newenvironment{history}{}{}
\usepackage{cleveref}
\crefname{originalproposition}{original proposition}{original propositions}
\Crefname{originalproposition}{Original proposition}{Original propositions}
\makeatletter
\newlabel{proposition-exists-minimal-model}{{0C2W}{source}{}{sourcepointer-1}{}}
\newlabel{proposition-exists-minimal-model@cref}{{[originalproposition][0][]0C2W}{[1][0][]source}{}{}{}}
\makeatother
\crefname{originallemma}{original lemma}{original lemmas}
\Crefname{originallemma}{Original lemma}{Original lemmas}
\makeatletter
\newlabel{lemma-exceptional-curves-dont-meet}{{0C6A}{source}{}{sourcepointer-2}{}}
\newlabel{lemma-exceptional-curves-dont-meet@cref}{{[originallemma][0][]0C6A}{[1][0][]source}{}{}{}}
\makeatother
\crefname{originallemma}{original lemma}{original lemmas}
\Crefname{originallemma}{Original lemma}{Original lemmas}
\makeatletter
\newlabel{resolve-lemma-factor-through-contraction}{{0C5J}{source}{}{sourcepointer-3}{}}
\newlabel{resolve-lemma-factor-through-contraction@cref}{{[originallemma][0][]0C5J}{[1][0][]source}{}{}{}}
\makeatother
\crefname{originallemma}{original lemma}{original lemmas}
\Crefname{originallemma}{Original lemma}{Original lemmas}
\makeatletter
\newlabel{resolve-lemma-proper-birational-regular-surfaces}{{0C5R}{source}{}{sourcepointer-4}{}}
\newlabel{resolve-lemma-proper-birational-regular-surfaces@cref}{{[originallemma][0][]0C5R}{[1][0][]source}{}{}{}}
\makeatother
\crefname{originallemma}{original lemma}{original lemmas}
\Crefname{originallemma}{Original lemma}{Original lemmas}
\makeatletter
\newlabel{resolve-lemma-birational-regular-surfaces}{{0C5S}{source}{}{sourcepointer-5}{}}
\newlabel{resolve-lemma-birational-regular-surfaces@cref}{{[originallemma][0][]0C5S}{[1][0][]source}{}{}{}}
\makeatother
\title{A curve of positive genus has a unique minimal model}
\author{The Stacks Project Authors}
\date{stacks-project:models.tex:local-snapshot-20261012; GNU FDL 1.2 or later}
\begin{document}
\maketitle
\noindent\textit{Editorial note.} The counted unit is the Stacks Project Lemma lemma-minimal-model-unique at source lines 4995-5000 together with its proof at source lines 5002-5048, a straight adjacent pairing in the chapter file. The chapter is the upstream file models.tex; the pool directory carries the longer name stacks-models and is not used as the chapter name anywhere in this card.\par\medskip
\noindent\textit{Editorial note.} This card has no context span: the statement and proof are self-contained apart from the five pointers, and the notions they use are recorded in the standard prerequisites.\par\medskip
\noindent\textit{Editorial note.} The retained proof names five results: the chapter's proposition on the existence of a minimal model, the chapter's lemma that exceptional curves of the first kind on a regular surface do not meet, and three results of the Resolution of Surfaces chapter. All five are declared as non-load-bearing navigation pointers with their stable Stacks tags looked up in the shared tag index; the first and the second belong to cards of this lane and their intervals are not reproduced here.\par\medskip
\noindent\textit{Editorial note.} The retained proof span starts at line 5002, the line of the opening marker read from the source, and contains exactly one proof environment.\par\medskip
\noindent\textit{Editorial note.} No citation occurs in the retained spans, so no bibliography entry is supplied.\par\medskip
\noindent\textit{Allowed standard prerequisites (editorial).}\begin{itemize}
\item Minimal models of a curve, exceptional curves of the first kind and the birational classification of regular surfaces. A model of a smooth projective curve over the fraction field of a discrete valuation ring is a flat proper scheme over the ring with that curve as generic fibre, and it is minimal when it admits no contraction to another model; the chapter proves that a minimal model always exists. Two exceptional curves of the first kind on a regular surface cannot meet, and a birational morphism between regular surfaces factors through the contraction of a finite sequence of such curves, a statement whose proof uses the classification of birational maps between regular surfaces. The author's own lemmas of the chapter and the three results of the Resolution of Surfaces chapter that the proof quotes are the allowed inputs.. The counted lemma is the uniqueness statement that the chapter's chapter-level results on good reduction depend on, and it is quoted by the lemma on the existence of a smooth model.
\end{itemize}\medskip
\ifdefined\hypertarget\hypertarget{sourcepointer-1}{}\fi
\noindent\textit{Source reference.} Original proposition 0C2W: focused-seeds/AGP/sources/stacks-snapshot-20261011/tags.txt, source lines 15171--15171. Navigation pointer invoked by the retained proof for the existence of a minimal model, which the comparison of two minimal models presupposes; that proposition is counted by a separate card of this lane, so its proof is not reproduced here. The number is the stable Stacks tag of the target, looked up in the shared tag index.\par\medskip
\ifdefined\hypertarget\hypertarget{sourcepointer-2}{}\fi
\noindent\textit{Source reference.} Original lemma 0C6A: focused-seeds/AGP/sources/stacks-snapshot-20261011/tags.txt, source lines 15290--15290. Navigation pointer invoked by the retained proof to keep the contractions of exceptional curves independent of one another; the cited statement is an allowed input of the same chapter and its proof is not reproduced here. The number is the stable Stacks tag of the target, looked up in the shared tag index.\par\medskip
\ifdefined\hypertarget\hypertarget{sourcepointer-3}{}\fi
\noindent\textit{Source reference.} Original lemma 0C5J: focused-seeds/AGP/sources/stacks-snapshot-20261011/tags.txt, source lines 15264--15264. Allowed input of the Resolution of Surfaces chapter quoted by the retained proof for the factorisation of a birational morphism between regular surfaces into contractions of exceptional curves; the pointer is registered so the reference resolves and the cited statement is not reproduced inside the span. The number is the stable Stacks tag of the target, looked up in the shared tag index.\par\medskip
\ifdefined\hypertarget\hypertarget{sourcepointer-4}{}\fi
\noindent\textit{Source reference.} Original lemma 0C5R: focused-seeds/AGP/sources/stacks-snapshot-20261011/tags.txt, source lines 15271--15271. Allowed input of the Resolution of Surfaces chapter quoted by the retained proof for the structure of a proper birational morphism between regular surfaces; the pointer is registered so the reference resolves and the cited statement is not reproduced inside the span. The number is the stable Stacks tag of the target, looked up in the shared tag index.\par\medskip
\ifdefined\hypertarget\hypertarget{sourcepointer-5}{}\fi
\noindent\textit{Source reference.} Original lemma 0C5S: focused-seeds/AGP/sources/stacks-snapshot-20261011/tags.txt, source lines 15272--15272. Allowed input of the Resolution of Surfaces chapter quoted by the retained proof for the classification of birational maps between regular surfaces; the pointer is registered so the reference resolves and the cited statement is not reproduced inside the span. The number is the stable Stacks tag of the target, looked up in the shared tag index.\par\medskip
\begin{lemma}
\label{lemma-minimal-model-unique}
Let $C$ be a smooth projective curve over $K$ with
$H^0(C, \mathcal{O}_C) = K$ and genus $> 0$.
There is a unique minimal model for $C$.
\end{lemma}

\begin{proof}
We have already proven the hard part of the lemma which is the existence
of a minimal model (whose proof relies on
resolution of surface singularities), see
Proposition \ref{proposition-exists-minimal-model}.
To prove uniqueness, suppose that $X$ and $Y$ are two
minimal models. By
Resolution of Surfaces, Lemma \ref{resolve-lemma-birational-regular-surfaces}
there exists a diagram of $S$-morphisms
$$
X = X_0 \leftarrow X_1 \leftarrow \ldots \leftarrow X_n = Y_m
\to \ldots \to Y_1 \to Y_0 = Y
$$
where each morphism is a blowup in a closed point. The
exceptional fibre of the morphism $X_n \to X_{n - 1}$ is an
exceptional curve of the first kind $E$. We claim that $E$ is
contracted to a point under the morphism $X_n = Y_m \to Y$.
If this is true, then $X_n \to Y$ factors through $X_{n - 1}$ by
Resolution of Surfaces, Lemma \ref{resolve-lemma-factor-through-contraction}.
In this case the morphism $X_{n - 1} \to Y$ is still a sequence of
contractions of exceptional curves by
Resolution of Surfaces, Lemma
\ref{resolve-lemma-proper-birational-regular-surfaces}.
Hence by induction on $n$ we conclude. (The base case $n = 0$ means
that there is a sequence of contractions
$X = Y_m \to \ldots \to Y_1 \to Y_0 = Y$
ending with $Y$. However as $X$ is a minimal model it contains
no exceptional curves of the first kind, hence $m = 0$ and $X = Y$.)

\medskip\noindent
Proof of the claim. We will show by induction on $m$ that any exceptional
curve of the first kind $E \subset Y_m$ is mapped to a point
by the morphism $Y_m \to Y$. If $m = 0$ this is clear because
$Y$ is a minimal model. If $m > 0$, then either
$Y_m \to Y_{m - 1}$ contracts $E$ (and we're done) or
the exceptional fibre $E' \subset Y_m$ of $Y_m \to Y_{m - 1}$
is a second exceptional curve of the first kind.
Since both $E$ and $E'$ are irreducible components of the special
fibre and since $g_C > 0$ by assumption, we conclude that
$E \cap E' = \emptyset$ by
Lemma \ref{lemma-exceptional-curves-dont-meet}.
Then the image of $E$ in $Y_{m - 1}$ is an exceptional
curve of the first kind (this is clear because the morphism
$Y_m \to Y_{m - 1}$ is an isomorphism in a neighbourhood of $E$).
By induction we see that $Y_{m - 1} \to Y$ contracts this curve
and the proof is complete.
\end{proof}

\end{document}
